% ============================================================
%  CARÁTULA  (Instrucción 3.3 + modelo del Anexo)
%  ------------------------------------------------------------
%  Reglas implementadas literalmente:
%   - Línea 1: UNIVERSIDAD MAYOR DE SAN SIMÓN  (Arial 16, negrilla)
%   - Línea 2: FACULTAD DE CIENCIAS Y TECNOLOGÍA (Arial 14, negrilla)
%   - Línea 3: DIRECCIÓN DE POSGRADO           (Arial 12, negrilla)
%     ...las tres en INTERLINEADO SIMPLE.
%   - Logo UMSS y logo FCyT: justificación IZQUIERDA.
%   - Logo Dirección de Posgrado: sección DERECHA, alineado con
%     la parte superior de la primera línea.
%   - TÍTULO: mayúsculas, Arial 20 negrilla, centrado; si ocupa
%     más de una línea, en forma de PIRÁMIDE INVERTIDA.
%   - 3 líneas abajo: "MONOGRAFÍA PRESENTADA..." (Arial 12, neg.)
%   - 5 líneas abajo: "POSTULANTE" (Arial 14, neg.) y a 1 cm
%     a la derecha ": NOMBRE".
%   - Al pie, centrado: "Cochabamba – Bolivia" y el AÑO (Arial 12).
%  La carátula NO muestra número de página pero SÍ cuenta (3.1).
% ============================================================

\begin{titlingpage}
\thispagestyle{empty}     % sin folio visible (pero la página se cuenta)
\interlineadoSimple       % toda la carátula va en interlineado simple

% ---------- Bloque superior: logos + encabezado institucional
% Se usa una fila de tres cajas superiores (\raisebox/[t]) para
% que los logos queden alineados con la PARTE SUPERIOR de la
% primera línea de texto, tal como pide la guía.
\noindent
\begin{minipage}[t]{0.21\textwidth}
  \vspace{0pt}
  \includegraphics[height=1.5cm]{imagenes/logo_umss.png}%
  \hspace{1mm}%
  \includegraphics[height=1.5cm]{imagenes/logo_facultad.png}
\end{minipage}%
\hfill
\begin{minipage}[t]{0.65\textwidth}
  \vspace{0pt}
  \centering
  {\fontsize{16}{18}\selectfont\bfseries UNIVERSIDAD MAYOR DE SAN SIMÓN\par}
  {\fontsize{14}{16}\selectfont\bfseries FACULTAD DE CIENCIAS Y TECNOLOGÍA\par}
  {\fontsize{12}{14}\selectfont\bfseries DIRECCIÓN DE POSGRADO\par}
\end{minipage}%
\hfill
\begin{minipage}[t]{0.13\textwidth}
  \vspace{0pt}
  \raggedleft
  \includegraphics[height=1.5cm]{imagenes/logo_posgrado.png}
\end{minipage}

\vspace{4.5cm}

% ---------- TÍTULO DE LA MONOGRAFÍA -------------------------
% \tituloMonografia se define en main.tex (datos del proyecto).
% Para la "pirámide invertida" cuando el título ocupa varias
% líneas, escribir el título con \\ donde se desee el corte:
% la primera línea más larga y las siguientes más cortas.
\begin{center}
  {\fontsize{20}{26}\selectfont\bfseries
   \MakeUppercase{\tituloMonografia}\par}
\end{center}

\vspace{3\baselineskip}   % "tres líneas por debajo"

% ---------- Leyenda de presentación -------------------------
\noindent
\begin{center}
  {\fontsize{12}{16}\selectfont\bfseries
   MONOGRAFÍA PRESENTADA PARA OBTENER EL CERTIFICADO\\[2pt]
   DE DIPLOMADO EN \MakeUppercase{\nombreDiplomado}\ \nombreVersion\ VERSIÓN\par}
\end{center}

\vspace{5\baselineskip}   % "cinco líneas por debajo"

% ---------- Postulante --------------------------------------
% "POSTULANTE" y, 1 cm a la derecha, ": NOMBRE DEL POSTULANTE".
\noindent\hspace{1.5cm}%
\begin{minipage}{\dimexpr\textwidth-1.5cm\relax}
  \raggedright   % sin justificar: evita que el nombre se estire entre palabras
  {\fontsize{14}{16}\selectfont\bfseries
   POSTULANTE\hspace{1cm}: \MakeUppercase{\nombrePostulante}\par}
\end{minipage}

\vfill

% ---------- Pie: lugar y año --------------------------------
\begin{center}
  {\fontsize{12}{14}\selectfont\bfseries Cochabamba -- Bolivia\par}
  {\fontsize{12}{14}\selectfont\bfseries \anioMonografia\par}
\end{center}

\end{titlingpage}
